% =====================================================================
%  OWNER: F. Only F edits this file; others use Overleaf comments.
%  Content: Industry pain points, cloud adoption motivation, KYC verification rules.
% =====================================================================
\section{Business Problem Framing and Requirements Analysis}\label{sec:business}

\subsection{Industry Background and Traditional Pain Points}\label{sec:business:background}

Customer onboarding is the first operational contact a bank has with a new
customer, and it is also the point at which regulatory risk and commercial
pressure collide. Under anti-money-laundering rules a bank must verify that an
applicant is who they claim to be before an account is opened: the identity
document has to be genuine and unexpired, and the person submitting it has to be
the person depicted on it. In a branch this is a manual task; in digital
onboarding it must be automated at scale, and the volume is strongly \emph{bursty}
because acquisition is driven by promotional campaigns rather than by a stable
daily arrival rate.

Traditional on-premises provisioning handles this burstiness badly, for three
reasons.

\begin{itemize}
  \item \textbf{Idle capacity versus over-provisioning.} Normal operating load
        typically consumes a small fraction of the provisioned server capacity,
        so most of the hardware is idle for most of the year. Provisioning
        additionally for the peak of a marketing campaign means paying for
        capital equipment that is used for a few days at a time.
  \item \textbf{Latency blow-ups at the peak.} When a campaign lands, a fixed
        infrastructure saturates: verification jobs queue behind each other, the
        queueing delay dominates the end-to-end time, and applicants abandon the
        flow. Industry data cited in the project proposal puts application
        abandonment at roughly 10\% when onboarding latency grows.
  \item \textbf{CapEx and operational drag.} Servers, storage, licences and the
        staff who maintain them are committed before a single application
        arrives. Every capacity decision is a multi-month procurement decision,
        which is the opposite of what a two-week campaign requires.
\end{itemize}

These are not tuning problems; they follow from the shape of the workload. A
verification pipeline is idle most of the time and extremely busy for short
intervals, which is precisely the profile that fixed-capacity infrastructure
prices worst.

\subsection{Cloud Adoption Motivation and Economic Value}\label{sec:business:cloud}

Cloud computing addresses the capacity-versus-cost conflict through two of the
essential characteristics defined by NIST~\cite{mell2011nist}, both of which this
project relies on directly:

\begin{itemize}
  \item \textbf{Rapid elasticity.} Compute is allocated per request rather than
        pre-purchased. When a burst of applicants arrives, an equal number of
        serverless execution environments are created in parallel and released
        as soon as each verification finishes, so the queue does not accumulate
        and the peak does not have to be provisioned for in advance.
  \item \textbf{Measured service (CapEx to OpEx).} Billing follows actual
        invocation count, execution duration and bytes stored. In the idle
        periods between campaigns the compute charge is zero, which converts a
        fixed capital commitment into a variable operating expense that tracks
        revenue.
\end{itemize}

The architectural consequence is that the unit of deployment stops being a
server and becomes an event. VeriCloud is therefore designed as an event-driven,
serverless pipeline in which every stage scales independently, following the
serverless application guidance of the AWS Well-Architected
Framework~\cite{aws2022serverlesslens}. The economic consequence is quantified in
Section~\ref{sec:cost}.

\subsection{KYC Compliance and Verification Rules}\label{sec:business:rules}

Automation is only acceptable if it remains auditable. The pipeline therefore
enforces a small set of quantitative rules, evaluated in a fixed order, so that
every application receives exactly one reproducible verdict. Table~\ref{tab:rules}
states them.

\begin{table}[htbp]
  \centering
  \small
  \caption{KYC decision rules, evaluated in order (first match wins).}
  \label{tab:rules}
  \begin{tabular}{@{}clll@{}}
    \toprule
    \# & Condition & Verdict & Reason code \\
    \midrule
    1 & OCR could not extract the required fields & \texttt{FAIL} & \texttt{ocr\_failed} \\
    2 & Document expired (\texttt{id\_valid = false}) & \texttt{FAIL} & \texttt{id\_expired} \\
    3 & Face comparison raised an error & \texttt{NEEDS\_HUMAN\_REVIEW} & \texttt{face\_error} \\
    4 & Similarity below the 80\% threshold & \texttt{NEEDS\_HUMAN\_REVIEW} & \texttt{similarity\_below\_80} \\
    5 & None of the above & \texttt{PASS} & \texttt{auto\_approved} \\
    \bottomrule
  \end{tabular}
\end{table}

Three design choices in Table~\ref{tab:rules} deserve comment. First, the rules
are \emph{ordered}: a missing extraction must never fall through to an approval,
and an expired document fails regardless of how well the faces match. Second, the
80\% similarity threshold is the only tunable in the automated path; it
is deliberately conservative, because the cost of a false acceptance in KYC is
far higher than the cost of a manual review. Third, anything that cannot be
decided automatically is routed to \texttt{NEEDS\_HUMAN\_REVIEW} rather than
being rejected, and the successfully extracted document fields are retained so
that the compliance officer does not have to re-key them. The system degrades to
a human, not to an error.

Files that are corrupt, empty or of an unsupported format are handled even
earlier: they fail validation before the workflow starts, are retried three
times, and are then quarantined in a dead-letter queue for 14 days. Such
applications never reach the decision table and never create a record, which
keeps the audit table free of unverifiable submissions.
